\chapter{Development Plan}
\label{chap:Chapter4}

This chapter presents the development plan for the proposed Boost-to-Rust
migration methodology. It details the planned reimplementation of selected
components in Rust, the design of stable interfaces, and the integration of
bidirectional \gls{FFI}. The plan is structured to
support incremental adoption while maintaining compliance with automotive
safety standards such as ISO-26262 and architectural principles from
AUTOSAR Adaptive.

The development activities are scheduled between February and July and are
organized to progressively reduce technical risk while preserving functional
continuity of the existing C++ system.

\section{Objectives and Scope}

The primary objectives of this development phase are:

\begin{itemize}
    \item To reimplement selected Boost-based components in Rust with equivalent
    functionality and improved memory-safety guarantees.
    \item To design explicit, minimal, and auditable interfaces between C++ and
    Rust components.
    \item To integrate Rust modules incrementally using bidirectional \gls{FFI},
    preserving deterministic behavior and freedom from interference.
    \item To validate the resulting hybrid system against safety, correctness,
    and performance expectations relevant to automotive software.
\end{itemize}

The scope of this chapter is limited to architectural and integration-level
development. Functional feature expansion beyond parity with the existing C++
implementation is explicitly excluded.

\section{Development Phases Overview}

The development plan is divided into six sequential phases:

\begin{enumerate}
    \item System analysis and Boost component selection
    \item Architectural refactoring and interface definition
    \item Rust reimplementation of selected components
    \item Bidirectional \gls{FFI} integration
    \item Validation and safety analysis
    \item Consolidation and documentation
\end{enumerate}

Each phase produces concrete artifacts that serve as inputs for the subsequent
phase, enabling traceability and controlled progression consistent with
ISO-26262 development practices.

\section{System Analysis and Component Selection}

\subsection*{February}

The initial phase focuses on analyzing the existing C++ system and identifying
Boost libraries that are relevant to automotive software architectures such as
those used in Eclipse SCORE–based environments.

The analysis includes:
\begin{itemize}
    \item Identification of Boost libraries in use and their functional roles
    \item Classification of components by criticality and \gls{ASIL} level
    \item Mapping of Boost abstractions to potential Rust-native equivalents
\end{itemize}

This phase establishes the baseline architecture and defines the migration
targets for subsequent development.

\section{Architectural Refactoring and Interface Design}

\subsection*{March}

In this phase, the selected C++ components are refactored to reduce semantic
ambiguity and prepare them for interoperability. Boost abstractions are used to
replace raw pointers, ad hoc ownership conventions, and implicit lifetimes.

Key activities include:
\begin{itemize}
    \item Encapsulation of migrated components behind stable interfaces
    \item Elimination of undefined behavior patterns
    \item Definition of ownership and lifetime semantics at component boundaries
\end{itemize}

This refactoring step does not change observable behavior but improves
analyzability and safety argumentation.

\section{Rust Reimplementation}

\subsection*{April}

During this phase, selected Boost-based components are reimplemented in Rust.
The Rust implementations are designed to be ownership-complete and memory-safe
by construction.

Design constraints include:
\begin{itemize}
    \item Explicit ownership and borrowing semantics
    \item Deterministic execution suitable for timing analysis
    \item Minimal use of \texttt{unsafe} code, isolated and documented
\end{itemize}

Functional equivalence with the original C++ components is verified through
unit testing and interface-level validation.

\section{Bidirectional \gls{FFI} Integration}

\subsection*{May}

This phase introduces bidirectional \gls{FFI} integration between C++ and Rust.
Stable C-compatible interfaces are defined, and responsibility for memory
allocation, ownership transfer, and error handling is made explicit.

Key concerns addressed include:
\begin{itemize}
    \item Isolation of unsafe code at \gls{FFI} boundaries
    \item Prevention of exception and panic propagation across language borders
    \item Alignment with AUTOSAR Adaptive service-oriented communication models
\end{itemize}

The \gls{FFI} boundary is treated as a safety-relevant interface and is subject to
additional review.

\section{Validation and Safety Analysis}

\subsection*{June}

Validation activities focus on ensuring that the hybrid C++/Rust system meets
functional, safety, and performance expectations.

This includes:
\begin{itemize}
    \item Functional regression testing
    \item Concurrency and race-condition analysis
    \item Review of freedom from interference between components of different
    \gls{ASIL} levels
\end{itemize}

The results of this phase support the safety argumentation required for
ISO-26262 compliant systems.

\section{Consolidation and Documentation}

\subsection*{July}

The final phase consolidates the development results and prepares the system
for evaluation and reporting.

Activities include:
\begin{itemize}
    \item Final architectural documentation
    \item Evaluation of migration feasibility and limitations
    \item Documentation of lessons learned and future work
\end{itemize}

This phase ensures that the development process and outcomes are reproducible
and auditable.

\section{Development Timeline}

Table \ref{tab:development_gantt} presents a Gantt-style overview of
the development plan covering the February--September period. The
activities are intentionally overlapping to reflect the iterative
nature of the migration process.

\begin{table}[H]

\centering

\renewcommand{\arraystretch}{1.3}

\begin{tabular}{|l|c|c|c|c|c|c|c|c|}

\hline

\textbf{Activity / Month} &
\textbf{Feb} &
\textbf{Mar} &
\textbf{Apr} &
\textbf{May} &
\textbf{Jun} &
\textbf{Jul} &
\textbf{Aug} &
\textbf{Sep} \\ \hline

System Analysis and Boost Selection
    & X & X &   &   &   &   &   &   \\ \hline

Architectural Refactoring and Interface Design
    &   & X & X &   &   &   &   &   \\ \hline

Rust Reimplementation
    &   &   & X & X &   &   &   &   \\ \hline

Bidirectional \gls{FFI} Integration
    &   &   &   & X & X &   &   &   \\ \hline

Validation and Safety Analysis
    &   &   &   &   & X & X & X &   \\ \hline

Performance Evaluation and Optimization
    &   &   &   &   &   & X & X &   \\ \hline

Consolidation and Documentation
    &   &   &   &   &   &   & X & X \\ \hline

Final Evaluation and Reporting
    &   &   &   &   &   &   &   & X \\ \hline

\end{tabular}

\caption{Gantt-style development plan for the Boost-to-Rust migration}

\label{tab:development_gantt}

\end{table}
